\bookchapter{บทที่ 0: สวัสดี ก่อนเริ่มต้น}

ก่อนจะเริ่มสำรวจอย่างเป็นทางการ มาคุยเรื่องกฎพื้นฐานกันก่อน เพื่อให้แน่ใจว่าสิ่งที่พูดต่อไปคุณจะเข้าใจทั้งหมด อย่ารีบร้อน ลับมีดไม่ทำให้ตัดฟืนช้า

\booksection{หนังสือเล่มนี้ช่วยอะไรคุณได้บ้าง? (และทำอะไรไม่ได้)}

มันช่วยให้คุณเข้าใจอย่างเป็นระบบว่าโทรศัพท์ใช้งานอย่างไร เข้าใจเหตุผลเบื้องหลังฟังก์ชันต่างๆ และไม่ว่าคุณจะเป็นฝ่ายแอนดรอยด์หรือฝ่ายแอปเปิล แนวคิดส่วนใหญ่ใช้ได้ทั้งหมด มันจะพาคุณลงมือฝึก ให้คุณอัปเกรดจาก “ใช้เป็น” เป็น “เข้าใจการใช้”

มันไม่ใช่คู่มือของโทรศัพท์รุ่นใดรุ่นหนึ่ง ไม่ใช่ข่าวเทคโนโลยีที่ไล่ตามผลิตภัณฑ์ใหม่ และไม่ใช่คู่มือสอนให้ติดโทรศัพท์ พูดตรงๆ ถ้าคุณคาดหวัง “วิธีดูวิดีโอบนมือถือให้สนุกขึ้น” อาจจะผิดหวัง เราสนใจมากกว่าว่าจะทำอย่างไรให้โทรศัพท์ทำงานเพื่อคุณอย่างแท้จริง ไม่ใช่กลับกัน

\booksection{คำอธิบาย “รหัสลับ” ของเรา}

เพื่อให้อธิบายได้ชัดเจนขึ้น หนังสือจะใช้สัญลักษณ์เล็กๆ สัญลักษณ์เหล่านี้เหมือนรหัสลับที่เราเข้าใจกัน:

\begin{itemize}
  \item \textbf{[ปุ่มชื่อ]}: หมายถึงปุ่มเฉพาะบนหน้าจอหรือชื่อแอป เช่น [ตั้งค่า], [ตกลง]
  \item \textbf{>}: หมายถึงลำดับการดำเนินการ เช่น "แตะ [ตั้งค่า] > [เสียง]" คือแตะตั้งค่าก่อน แล้วแตะเสียง
  \item \textbf{ตัวหนา}: สำคัญมาก! อาจเป็นแนวคิดสำคัญหรือคำเตือนความปลอดภัย โปรดสังเกตให้ดี
  \item \emph{ตัวเอียง}: ความคิดเล็กๆ น้อยๆ เพิ่มเติมหรือการเปรียบเทียบ ช่วยให้คุณเข้าใจได้ดีขึ้น
\end{itemize}

\booksection{แอนดรอยด์และ iOS เราพูดถึงทั้งคู่}

คุณเดาสิ บางอย่างเรียกชื่อต่างกัน แต่ทำหน้าที่เหมือนกัน เช่น: แอปพลิเคชัน (แอนดรอยด์) และ App (iOS), ตั้งค่า (ทั้งคู่เรียกตั้งค่า), ไฟล์เมเนเจอร์ (แอนดรอยด์) และ ไฟล์ (iOS) รู้ว่าเป็นสิ่งเดียวกันก็พอ อย่าให้คำศัพท์ทำให้สับสน

นอกจากนี้ ในฝ่ายแอนดรอยด์ยังมีระบบปฏิบัติการฮาร์โมนีโอเอสของหัวเว่ย ซึ่งมีเอกลักษณ์เฉพาะตัว เราจะพูดถึงเป็นพิเศษในภายหลัง สรุปคือ ไม่ว่าคุณใช้โทรศัพท์อะไร หนังสือเล่มนี้ก็ช่วยคุณได้

\booksection{ความจุโทรศัพท์: ทำไมซื้อ 128G แต่ใช้ได้ไม่ถึง 110G?}

หน่วยความจุทุกคนรู้: 1TB = 1024GB, 1GB = 1024MB, 1MB = 1024KB... แต่มีข้อสังเกต: ผู้ผลิตฮาร์ดดิสก์มักคำนวณแบบฐาน 1000 ในขณะที่ระบบคำนวณแบบฐาน 1024 ดังนั้นฮาร์ดดิสก์ “128G” ที่คุณซื้อ ความจุจริงจะ “ลดลง”

แต่สาเหตุหลักที่ความจุโทรศัพท์ของคุณลดลงมี 3 ประการ:

\begin{enumerate}
  \item \textbf{ระบบปฏิบัติการใช้พื้นที่เอง}: ระบบปฏิบัติการเหมือนผนังรับน้ำหนักของบ้าน ต้องมีอยู่ ระบบปัจจุบันมักใช้พื้นที่ 10-20GB ซึ่งเป็นเรื่องปกติ
  \item \textbf{แอปที่ผู้ผลิตติดตั้งมาให้}: แอปที่มาพร้อมกับโทรศัพท์ใหม่ก็ใช้พื้นที่ บางตัวถอนการติดตั้งได้ บางตัวไม่ได้
  \item \textbf{การสูญเสียจากการแบ่งพาร์ทิชันและการฟอร์แมต}: เหตุผลทางเทคนิค เหมือนผ้าบุของกระเป๋าเป้ที่ก็ใช้พื้นที่
\end{enumerate}

\begin{tcbinfo}{คำแนะนำเล็กๆ ในการซื้อโทรศัพท์}
โทรศัพท์ 128GB ในปัจจุบันไม่พอใช้แล้ว 256GB เป็นจุดเริ่มต้นที่ปลอดภัย ถ้าคุณชอบถ่ายรูป เล่นเกม แนะนำให้ไปเลย 512GB หรือใหญ่กว่า จำไว้ พื้นที่เก็บข้อมูลนี่ ขอให้เหลือใช้ดีกว่าไม่พอ

\end{tcbinfo}

\begin{tcbpractice}{ลงมือ: ดู "บัตรประจำตัว" ของโทรศัพท์คุณ}
เปิด [ตั้งค่า] หาเกี่ยวกับโทรศัพท์ (หรือตัวเลือกที่คล้ายกัน) จดรุ่นโทรศัพท์ เวอร์ชันระบบ พื้นที่เก็บข้อมูลรวมและพื้นที่ว่าง นี่คือก้าวแรกในการทำความรู้จักมัน ดูสุขภาพแบตเตอรี่ด้วยเพื่อให้รู้สถานะ

\end{tcbpractice}

\booksection{โทรศัพท์ของคุณคือ “คลังสินทรัพย์ดิจิทัล”}

ก่อนเริ่มความรู้เฉพาะทาง มาสร้างความเข้าใจหลักกันก่อน:\textbf{โทรศัพท์ไม่ใช่แค่เครื่องมือสื่อสารอีกต่อไป มันคือคลังสินทรัพย์ดิจิทัลที่สำคัญที่สุดของเรา}รายชื่อผู้ติดต่อ ประวัติแชท รูปภาพครอบครัว เอกสารงาน แอปธนาคาร QR ชำระเงิน ข้อมูลสุขภาพ บัตรประจำตัวอิเล็กทรอนิกส์... ข้อมูลเหล่านี้รวมกันแล้ว มีค่ามากกว่าตัวเครื่องโทรศัพท์เสียอีก

\begin{tcbinfo}{ทำไมนี่จึงสำคัญ?}
ถ้าโทรศัพท์หายหรือพัง คุณจะไม่เพียงสูญเสียฮาร์ดแวร์มูลค่าไม่กี่พันบาท\textbf{แต่อาจสูญเสียข้อมูลที่สะสมมาหลายปีและไม่สามารถเรียกคืนได้}— รูปภาพการเติบโตของลูก เสียงของญาติที่ล่วงลับ สัญญาสำคัญที่ไม่ได้สำรอง การจัดระเบียบและนิสัยการป้องกันอย่างเป็นวิทยาศาสตร์ สามารถปกป้องคุณได้ใน 3 มิติ:

\begin{itemize}
  \item \textbf{เพิ่มประสิทธิภาพการทำงาน}: การจัดประเภทไฟล์ที่เป็นระบบทำให้ค้นหาเร็วขึ้น โทรศัพท์ไม่ค้างในระยะยาว
  \item \textbf{ยืดอายุการใช้งาน}: การจัดการพื้นที่เก็บข้อมูลที่เหมาะสมลดภาระการอ่าน/เขียนของชิป ทั้งแบตเตอรี่และพื้นที่เก็บข้อมูลทนทานขึ้น
  \item \textbf{รับประกันความปลอดภัยข้อมูล}: การจัดการแบบชั้น + สำรองข้อมูลเป็นประจำ + จำกัดสิทธิ์ ทำให้ "ภัยพิบัติข้อมูล" ห่างจากคุณ
\end{itemize}

\end{tcbinfo}

\booksubsection{ปรัชญาหลักของเรา: “ป้องกันเป็นหลัก ป้องกันและรักษาควบคู่กัน”}

เหมือนการดูแลสุขภาพ เราไม่สามารถรอจนป่วยแล้วค่อยไปหาหมอ\textbf{การจัดการโทรศัพท์ก็ต้องเริ่มจาก “นิสัยในชีวิตประจำวัน”}หนังสือเล่มนี้ไม่ได้เน้นที่เทคนิคขั้นสูงใดๆ แต่เป็นหลักการพื้นฐานเหล่านี้:

\begin{itemize}
  \item \textbf{ป้องกันเป็นหลัก}: ติดตั้งแอปผ่านร้านค้าทางการเท่านั้น ลบข้อมูลสำคัญในรูปก่อนส่ง ไม่ใช้รหัสผ่านซ้ำ — "การกระทำเล็กๆ" เหล่านี้มีต้นทุนต่ำมาก แต่ป้องกันความเสี่ยงส่วนใหญ่ได้
  \item \textbf{ตรวจสุขภาพเป็นประจำ}: จัดโฟลเดอร์ดาวน์โหลดทุกเดือน ตรวจสอบสิทธิ์ทุกไตรมาส สำรองข้อมูลทั้งหมดทุกครึ่งปี — ระยะห่างคงที่มีประสิทธิภาพมากกว่าดับไฟเฉพาะหน้าร้อยเท่า
  \item \textbf{เหลือทางถอย}: เปิดการซิงค์คลาวด์ จดรหัสผ่านบัญชี เขียน "รายการมรดกดิจิทัล" — รับประกันว่าเมื่อเกิดเหตุฉุกเฉิน คุณและครอบครัวสามารถรับมือได้อย่างสงบ
\end{itemize}
